\section{Методы условной оптимизации}

Теперь задача решается не на всём пространстве, а на допустимом множестве:
\[
  \varphi_0(x) \longrightarrow \min_{x \in S} ,
  \qquad S \subset \R^n .
\]

\begin{idea}[title={Общая картина: зачем вообще нужны отдельные методы}]
Все методы предыдущих глав строили последовательность спуска, свободно двигаясь
по $\R^n$. Ограничение $x \in S$ ломает эту схему: шаг вдоль антиградиента может
вывести из $S$, а минимум $\varphi_0$ на $S$, как правило, лежит на границе,
где $\nabla\varphi_0 \ne 0$ — то есть привычный критерий остановки
$\norm{\nabla \varphi_0} < \eps$ не сработает никогда.

Отсюда три принципиально разных подхода, и вся оставшаяся часть курса — про них:
\begin{enumerate}[label=\arabic*.,left=0pt,itemsep=3pt]
  \item \textbf{Методы возможных направлений} — остаться внутри $S$ и на каждом
        шаге искать направление, которое \emph{одновременно} не выводит из $S$ и
        уменьшает $\varphi_0$. Представитель — метод Зойтендейка
        (\S\ref{subsec:zoutendijk}).
  \item \textbf{Методы отсечения} — заменить $S$ множеством, описываемым более
        простыми условиями (обычно многогранником), решить задачу на нём и
        отсечь часть, где оптимума заведомо нет; повторять.
  \item \textbf{Методы штрафов} — вообще отказаться от ограничений, встроив их
        в целевую функцию, и решать последовательность задач \emph{без}
        ограничений уже известными методами.
\end{enumerate}
\end{idea}

\subsection{Методы возможных направлений}

\begin{definition}\label{def:feasible-direction}
Пусть $x \in S$, $\sigma \in \R^n$. Вектор $\sigma$ называется
\textbf{возможным направлением} в точке $x$, если
\[
  \exists\, \delta_0 > 0 : \qquad x + \delta\sigma \in S
  \quad \forall \delta \in [0, \delta_0] .
\]
\end{definition}

\inputfig{figures/fig-feasible-direction}

\begin{idea}
Во внутренней точке $S$ возможным является любое направление — отсюда и
совпадение с безусловной задачей вдали от границы. Всё содержание определения
проявляется на границе: там конус возможных направлений сужается, и метод
обязан выбирать направление спуска \emph{внутри} этого конуса. Пересечение двух
условий — «возможное» ($\sigma$ не выводит из $S$) и «убывающее»
($\nabla\T\varphi_0(x)\sigma < 0$) — и есть то, что ищет метод Зойтендейка.
\end{idea}

\subsection{Методы отсечения}

Аппроксимируем $S$ множеством, описываемым более простыми условиями. Построим
вспомогательную задачу, которая позволит отсечь от этого множества множество
точек, в котором оптимума заведомо нет; потом вторую, и так далее.

\inputfig{figures/fig-cutting-plane}

\begin{idea}
Схема работы: исходное $S$ заменяют объемлющим многогранником $P_0 \supset S$,
на котором задача решается просто (например, линейным программированием). Если
найденный оптимум $x^{(0)}$ попал в $S$ — он и есть решение. Если нет, через
него проводят \textbf{отсекающую плоскость}: гиперплоскость, которая отрезает
$x^{(0)}$, но не отрезает ни одной точки $S$. Получается $P_1 \subset P_0$,
по-прежнему содержащее $S$, и всё повторяется. Множества $P_k$ сжимаются к $S$
в окрестности оптимума.
\end{idea}

\subsection{Методы штрафов}

\paragraph{Направление 1 (внешний штраф).}
Строим последовательность спусков в области, но неявно: записываем
вспомогательную задачу, содержащую \emph{штраф за выход} за границу области, и
отказываемся от решения задачи с ограничениями. При решении вспомогательной
задачи значение функции цели в области совпадает со значением в области
исходной задачи — значит можно решать задачу без ограничений.

\paragraph{Направление 2 (барьерные функции, внутренний штраф).}
Строим вспомогательную задачу так, чтобы \emph{не выпускать} точку за границу,
и решаем задачу без ограничений.

\inputfig{figures/fig-penalty-barrier}

\begin{idea}
Формально обе конструкции выглядят так:
\[
  \text{внешний штраф:}\quad
  \Phi_r(x) = \varphi_0(x) + r \sum_{i=1}^{m} \big(\max\{0, \varphi_i(x)\}\big)^2 ,
  \qquad r \to +\infty ;
\]
\[
  \text{барьер:}\quad
  \Phi_r(x) = \varphi_0(x) - r \sum_{i=1}^{m} \frac{1}{\varphi_i(x)} ,
  \qquad r \to +0 .
\]
Разница принципиальная:
\begin{itemize}[nosep,left=0pt]
  \item \emph{внешний} штраф равен нулю внутри $S$ и растёт снаружи, поэтому
        внутри области $\Phi_r \equiv \varphi_0$ — это и есть «значение функции
        цели совпадает». Приближения идут \emph{снаружи} $S$ и допустимыми
        становятся лишь в пределе;
  \item \emph{барьер} обращается в $+\infty$ при подходе к границе изнутри, то
        есть точка физически не может выйти из $S$. Зато он искажает $\varphi_0$
        и внутри области, и это искажение снимается лишь при $r \to 0$.
\end{itemize}
Платят за удобство обусловленностью: при $r \to \infty$ (или $r \to 0$)
вспомогательная функция становится всё более овражной, и безусловные методы на
ней работают всё хуже. Поэтому решают не одну задачу, а последовательность с
постепенно меняющимся $r$, беря решение предыдущей как начальное приближение
следующей.
\end{idea}
